\documentclass[11pt]{article}
% ===== Préambule commun aux documents Séries Chronologiques =====
\usepackage[utf8]{inputenc}
\usepackage[T1]{fontenc}
\usepackage[french]{babel}
\usepackage{lmodern}
\usepackage{amsmath,amssymb,amsthm,mathtools}
\usepackage{geometry}
\geometry{a4paper,margin=2.3cm}
\usepackage{enumitem}
\usepackage{array,booktabs,longtable}
\usepackage{xcolor}
\usepackage[most]{tcolorbox}
\usepackage{etoolbox}
\usepackage{graphicx}
\graphicspath{{figs/}}
\usepackage{caption}
\captionsetup{font=small,labelfont={bf,color=insea}}
\usepackage{fancyhdr}
\usepackage{titlesec}
\usepackage{hyperref}
\hypersetup{colorlinks=true,linkcolor=blue!55!black,urlcolor=blue!55!black,
  citecolor=blue!55!black,bookmarksopen=true}
\usepackage{bookmark}

% --- Couleurs ---
\definecolor{insea}{RGB}{0,76,128}
\definecolor{insealight}{RGB}{232,240,247}
\definecolor{vert}{RGB}{0,120,70}
\definecolor{vertlight}{RGB}{232,245,238}
\definecolor{orange}{RGB}{190,90,0}
\definecolor{orangelight}{RGB}{252,242,228}
\definecolor{gris}{RGB}{90,90,90}

% --- En-têtes / pieds ---
\pagestyle{fancy}
\fancyhf{}
\fancyhead[L]{\small\color{gris}\leftmark}
\fancyhead[R]{\small\color{gris}Séries Chronologiques}
\fancyfoot[C]{\thepage}
\renewcommand{\headrulewidth}{0.4pt}
\renewcommand{\footrulewidth}{0pt}

% --- Titres de section colorés ---
\titleformat{\section}{\Large\bfseries\color{insea}}{\thesection.}{0.6em}{}
\titleformat{\subsection}{\large\bfseries\color{insea!85!black}}{\thesubsection.}{0.6em}{}
\titleformat{\subsubsection}{\bfseries\color{gris}}{\thesubsubsection.}{0.6em}{}

% --- Théorèmes (pour le cours) ---
\theoremstyle{definition}
\newtheorem{definition}{Définition}[section]
\newtheorem{exemple}{Exemple}[section]
\newtheorem{remarque}{Remarque}[section]
\newtheorem{methode}{Méthode}[section]
\theoremstyle{plain}
\newtheorem{theoreme}{Théorème}[section]
\newtheorem{proposition}{Proposition}[section]
\newtheorem{corollaire}{Corollaire}[section]
\newtheorem{propriete}{Propriété}[section]

% --- Boîtes ---
\newtcolorbox{recette}[1][]{colback=vertlight,colframe=vert,
  fonttitle=\bfseries,title={Recette d'examen\ifblank{#1}{}{ --- #1}}}
\newtcolorbox{attention}[1][]{colback=orangelight,colframe=orange,
  fonttitle=\bfseries,title={Attention\ifblank{#1}{}{ --- #1}}}
\newtcolorbox{rappel}[1][]{colback=insealight,colframe=insea,
  fonttitle=\bfseries,title={Rappel de cours\ifblank{#1}{}{ --- #1}}}
\newtcolorbox{reponse}[1][]{colback=vertlight,colframe=vert!70!black,
  left=2mm,right=2mm,top=1mm,bottom=1mm,boxrule=0.8pt,#1}

% --- Notations Séries Chronologiques ---
\newcommand{\E}{\mathbb{E}}
\newcommand{\Var}{\operatorname{Var}}
\newcommand{\Cov}{\operatorname{Cov}}
\newcommand{\Corr}{\operatorname{Corr}}
\newcommand{\cor}{\operatorname{cor}}
\newcommand{\BB}{\operatorname{BB}}
\newcommand{\R}{\mathbb{R}}
\newcommand{\Z}{\mathbb{Z}}
\newcommand{\N}{\mathbb{N}}
\newcommand{\Xhat}{\widehat{X}}
\newcommand{\phihat}{\widehat{\phi}}
\newcommand{\rhohat}{\widehat{\rho}}
\newcommand{\gammahat}{\widehat{\gamma}}
\newcommand{\eps}{\varepsilon}
\DeclareMathOperator{\sgn}{sgn}

% Étiquette « bonne réponse » pour les QCM (utilisée dans le corrigé)
\newcommand{\bonne}[1]{\textbf{\color{vert}#1}}
\newcommand{\horsprog}{\textsuperscript{\color{orange}\,$\star$}}


\title{\vspace{-1cm}\textbf{\color{insea}Séries Chronologiques}\\[2mm]
\large Corrigé détaillé des questions et exercices\\[1mm]
\normalsize (solutions fondées sur le cours de F.~BADAOUI, INSEA)}
\author{Préparation au module}
\date{\today}

\begin{document}
\maketitle
\thispagestyle{empty}

\begin{tcolorbox}[colback=insealight,colframe=insea,title={Mode d'emploi}]
Ce document corrige \textbf{toutes} les questions du recueil \emph{Questions\_SC.pdf}, en
conservant la \textbf{même numérotation} (Q\,A.1, Q\,B.2, \dots). Chaque corrigé est
\textbf{justifié à partir du cours} (notations de F.~BADAOUI). Les notions \horsprog{}
dépassant le cours sont traitées dans la dernière section.

\smallskip\noindent\textbf{Conventions de cours rappelées.}
Opérateur retard $B$ : $B^kX_t=X_{t-k}$.
ARMA$(p,q)$ : $\Phi_p(B)X_t=\mu+\Theta_q(B)a_t$ avec
$\Phi_p(B)=1-\phi_1B-\dots-\phi_pB^p$, $\Theta_q(B)=1-\theta_1B-\dots-\theta_qB^q$,
$a_t\sim\BB(0,\sigma^2)$.
\textbf{Stationnarité} (causalité) d'un AR : racines de $\Phi_p(B)=0$ de module $>1$.
\textbf{Inversibilité} d'un MA : racines de $\Theta_q(B)=0$ de module $>1$.
\end{tcolorbox}

\tableofcontents
\newpage

% =====================================================================
\section{QCM théoriques et questions de cours}
\markboth{QCM théoriques}{}

\subsection*{Concepts de base et décomposition}

\textbf{Q\,A.1 --- Processus stationnaire.}
\begin{reponse}Réponse : \bonne{A}.\end{reponse}
Un processus SSL a une espérance \emph{constante} et une autocovariance indépendante de $t$.
Une tendance fait varier $\E(X_t)$ avec $t$ ; une saisonnalité crée une dépendance
périodique : dans les deux cas la condition de stationnarité est violée. Un processus
stationnaire modélise donc des séries \emph{sans} tendance ni saisonnalité (c'est pourquoi on
\og stationnarise \fg{} d'abord par différenciation/désaisonnalisation).

\textbf{Q\,A.2 --- Moyenne mobile.}
\begin{reponse}Réponse : \bonne{B} (filtre linéaire, pas forcément centré ou symétrique).\end{reponse}
Une moyenne mobile est un \emph{filtre linéaire} $M(X)_t=\sum_{i=-m_1}^{m_2}\lambda_i X_{t-i}$
(combinaison linéaire d'observations). Elle n'est ni nécessairement centrée ($m_1=m_2$) ni
symétrique ($\lambda_{-i}=\lambda_i$) ni normalisée ($\sum\lambda_i=1$) : ce sont des
propriétés \emph{supplémentaires} souhaitables, non constitutives de la définition.

\textbf{Q\,A.3 --- Moyenne mobile MM4.}
\begin{reponse}Réponse : \bonne{C} (conserve les tendances linéaires, absorbe les
saisonnalités de période 4).\end{reponse}
Une moyenne mobile d'ordre $p$ (centrée) \emph{laisse invariante} une tendance affine et
\emph{annule} (absorbe) une composante saisonnière de période $p$. La MM4 (centrée :
$\frac12$ aux extrémités) élimine donc la saisonnalité \emph{trimestrielle} (période 4) tout
en préservant la tendance linéaire ; on l'utilise précisément pour extraire la tendance d'une
série trimestrielle.

\textbf{Q\,A.4 --- Algorithme X11.}
\begin{reponse}Réponses correctes : \bonne{« contient deux boucles pour évaluer la série
CVS »} et \bonne{« intègre conjointement régression linéaire et moyenne mobile »}.\end{reponse}
X11 procède de façon itérative (estimation tendance par MM $\to$ coefficients saisonniers
$\to$ correction $\to$ ré-estimation) en combinant régressions et moyennes mobiles.
L'affirmation \og l'algorithme le plus sophistiqué connu \fg{} est \emph{fausse}
(il existe X-12-ARIMA, X-13ARIMA-SEATS).

\textbf{Q\,A.5 --- Décomposition $=$ enlever la saisonnalité ?}
\begin{reponse}Réponse : \bonne{Faux}.\end{reponse}
La décomposition vise à \emph{isoler les trois composantes} $T_t$ (tendance), $S_t$
(saisonnalité) et $E_t$ (résidu). La désaisonnalisation (série CVS) n'en est qu'une
\emph{application} parmi d'autres (étude de la tendance, prévision, \dots).

\textbf{Q\,A.6 --- Première moyenne mobile.}
\begin{reponse}Réponse : \bonne{Vrai}.\end{reponse}
La MM employée pour extraire la tendance doit \emph{conserver} la tendance, \emph{absorber}
la saisonnalité et \emph{réduire} la variance du bruit (effet de lissage). C'est exactement le
rôle d'une MM d'ordre égal à la période.

\textbf{Q\,A.7 --- Lissage exponentiel simple.}
\begin{reponse}Réponse : \bonne{Vrai}.\end{reponse}
Le lissage exponentiel simple ajuste \emph{localement une constante} : il ne modélise ni
tendance ni saisonnalité. Pour la tendance on utilise le lissage double (Holt), pour la
saisonnalité le lissage de Holt--Winters.

\textbf{Q\,A.8 --- Holt--Winters.}
\begin{reponse}Réponses correctes : \bonne{A} et \bonne{C}.\end{reponse}
Holt--Winters généralise les lissages simple/double (donc plus complet/efficace) et prend en
compte \emph{à la fois} une tendance \emph{et} une saisonnalité. L'item B est faux.

\subsection*{Caractérisation AR / MA / ARMA / SARIMA}

\textbf{Q\,A.9 --- Processus AR.}
\begin{reponse}Réponse : \bonne{B}.\end{reponse}
Pour un AR$(p)$ : la FAC (ACF) \emph{décroît} (exponentiellement / en sinusoïde amortie) et la
FAP (PACF) \emph{s'annule} pour $k>p$ ($P_p=\phi_p\ne0$, $P_k=0$, $k>p$). C'est la dualité du
cours.

\textbf{Q\,A.10 --- Processus ARMA.}
\begin{reponse}Réponse : \bonne{D} (aucune des propriétés précédentes).\end{reponse}
Pour un ARMA$(p,q)$ \emph{ni} l'ACF \emph{ni} la PACF ne s'annulent à partir d'un rang : les
\emph{deux} décroissent (mélange des comportements AR et MA). Aucune coupure nette.

\textbf{Q\,A.11 --- Processus SARIMA.}
\begin{reponse}Réponse : \bonne{D} (tendance et/ou saisonnalité).\end{reponse}
Le SARIMA généralise l'ARIMA en ajoutant une partie saisonnière
$\Phi(B)\Phi_S(B^s)(1-B)^d(1-B^s)^D X_t=\Theta(B)\Theta_S(B^s)a_t$ : il traite à la fois
tendance (via $(1-B)^d$) et saisonnalité (via $(1-B^s)^D$). Il n'est pas stationnaire en
général (item A faux).

\textbf{Q\,A.12 --- Tout MA stationnaire possède une représentation AR.}
\begin{reponse}Réponse : \bonne{Faux}.\end{reponse}
Un MA$(q)$ est \emph{toujours} stationnaire, mais il n'admet une représentation
$\mathrm{AR}(\infty)$ que s'il est \emph{inversible} (racines de $\Theta_q(B)=0$ de module
$>1$). Contre-exemple : $X_t=a_t-a_{t-1}$, racine $B=1$, non inversible.

\textbf{Q\,A.13 --- Validation d'un modèle ARMA.}
\begin{reponse}Réponse essentielle : \bonne{B} (blancheur du résidu) ; en pratique on
vérifie aussi \bonne{A} (significativité des paramètres) et \bonne{C} (normalité des
résidus).\end{reponse}
Le c\oe ur de la validation est que le résidu estimé soit un \emph{bruit blanc} (test de
Ljung--Box / examen ACF-PACF des résidus). On contrôle en complément la significativité des
coefficients (tests de Student) et la normalité (Jarque--Bera) pour les intervalles de
prévision.

\textbf{Q\,A.14 --- Ljung--Box sur une série stationnaire.}
\begin{reponse}Réponse : \bonne{Faux}.\end{reponse}
\emph{Stationnaire} $\ne$ \emph{bruit blanc}. Un AR(1) stationnaire a des autocorrélations non
nulles : la statistique de Ljung--Box (dont $H_0$ : \og toutes les autocorrélations sont
nulles \fg) y serait \emph{rejetée}. Le test de Ljung--Box s'applique aux \emph{résidus} d'un
modèle ajusté, pas à la série brute.

\subsection*{Questions de cours}

\textbf{Q\,A.15 --- Deux types de stationnarité.}
\emph{Stationnarité stricte (forte)} : la loi de $(X_{t_1},\dots,X_{t_n})$ est invariante par
translation temporelle. \emph{Stationnarité au sens large (SSL, second ordre)} : $\E(X_t)$
constante, $\E(X_t^2)<\infty$ et $\gamma(k)=\Cov(X_t,X_{t+k})$ indépendante de $t$.
La SSL est \emph{plus pratique} car elle ne porte que sur les deux premiers moments
(vérifiable empiriquement), alors que la stationnarité stricte porte sur toute la loi.

\textbf{Q\,A.16 --- Strict $\Rightarrow$ second ordre.}
Si $(X_t)$ est strictement stationnaire \emph{et} de carré intégrable ($\E(X_t^2)<\infty$),
alors il est SSL : la stationnarité stricte impose l'invariance de toute la loi, donc en
particulier de $\E(X_t)$ et de $\gamma(k)$, dès que ces moments existent.

\textbf{Q\,A.17 --- Corrélogramme.}
Le corrélogramme (graphe de $\rhohat(k)$) permet d'\emph{identifier la structure de
dépendance} : détecter une non-stationnarité (décroissance lente $\approx1$), une saisonnalité
(pics aux multiples de la période), et l'ordre $q$ d'un MA (l'ACF s'annule après $q$).

\textbf{Q\,A.18 --- AR$(p)$ par la FAP.}
Un AR$(p)$ est caractérisé par une \emph{fonction d'autocorrélation partielle qui s'annule
pour $k>p$} : $P_k=0,\ \forall k>p$ et $P_p=\phi_p\ne0$. La FAP \og coupe \fg{} à l'ordre $p$.

\textbf{Q\,A.19 --- Deux tests de validation ARMA gaussien.}
(i) \emph{Test de Ljung--Box} sur les résidus : $H_0$ \og résidus = bruit blanc \fg, statistique
$\mathrm{LB}=n(n+2)\sum_{k=1}^{K}\frac{\rhohat_k^2}{n-k}\sim\chi^2_{K-p-q}$ sous $H_0$.
(ii) \emph{Test de normalité de Jarque--Bera} (ou test de Student sur la significativité des
paramètres). On peut aussi citer le test du Portmanteau (Box--Pierce).

\textbf{Q\,A.20 --- Qualité prédictive.}
Critères d'erreur de prévision \emph{hors échantillon} : RMSE $=\sqrt{\frac1h\sum(X_{n+i}-\Xhat_{n+i})^2}$,
MAE, MAPE ; ou critère in-sample comme le $R^2$ ajusté. (Les critères AIC/BIC servent surtout
à la sélection.)

\textbf{Q\,A.21 --- Deux critères de sélection AR$(p)$ gaussien.}
\emph{AIC} $=-2\ln L+2k$ et \emph{BIC} (Schwarz) $=-2\ln L+k\ln n$, où $k$ est le nombre de
paramètres. On retient le modèle qui \emph{minimise} le critère ; le BIC pénalise davantage la
complexité (modèles plus parcimonieux).

% =====================================================================
\section{Stationnarité, causalité et inversibilité}
\markboth{Stationnarité, causalité, inversibilité}{}

\textbf{Q\,B.1 --- $X_t=(-1)^tz$, $\E(z)=0$, $\Var(z)=1$.}
\[
\E(X_t)=(-1)^t\E(z)=0,\qquad
\gamma(k)=\Cov(X_t,X_{t+k})=(-1)^t(-1)^{t+k}\E(z^2)=(-1)^{2t+k}=(-1)^k .
\]
$\E(X_t)=0$ (constante) et $\gamma(k)=(-1)^k$ \emph{ne dépend pas de $t$}. Donc
\begin{reponse}\bonne{$X_t$ est stationnaire au sens large} (et même $\gamma(0)=1$,
$\rho(k)=(-1)^k$).\end{reponse}
C'est cohérent avec l'exercice du cours : si $X$ est SSL centré, $(-1)^tX_t$ l'est aussi.
\begin{remarque}
Ce processus n'est \emph{pas} ergodique (une seule réalisation de $z$ détermine toute la
trajectoire), mais la stationnarité SSL ne concerne que les deux premiers moments : elle est
bien vérifiée.
\end{remarque}

\textbf{Q\,B.2 --- $X_t=0.8X_{t-1}-0.15X_{t-2}+a_t-0.3a_{t-1}$.}
Forme : $\Phi(B)X_t=\Theta(B)a_t$ avec $\Phi(B)=1-0.8B+0.15B^2$, $\Theta(B)=1-0.3B$.
\emph{Stationnarité} (racines de $\Phi$) : $0.15B^2-0.8B+1=0\Rightarrow
B=\frac{0.8\pm\sqrt{0.64-0.6}}{0.3}=\frac{0.8\pm0.2}{0.3}\in\{2,\ 3.33\}$, modules $>1$
$\Rightarrow$ \emph{stationnaire}.
\emph{Inversibilité} (racine de $\Theta$) : $B=1/0.3=3.33>1\Rightarrow$ \emph{inversible}.
\begin{reponse}Réponse : \bonne{(a) stationnaire et inversible}.\end{reponse}

\textbf{Q\,B.3 --- $(1-1.66B+0.66B^2)X_t=(1-B)a_t$ : représentation AR d'ordre fini ?}
On factorise : $1-1.66B+0.66B^2=(1-B)(1-0.66B)$. L'équation devient
\[
(1-B)(1-0.66B)X_t=(1-B)a_t .
\]
Le facteur commun $(1-B)$ se simplifie, d'où $(1-0.66B)X_t=a_t$, soit un \emph{AR(1)} :
\begin{reponse}Réponse : \bonne{Vrai} --- après simplification du facteur commun, le processus
est un AR(1) $X_t=0.66X_{t-1}+a_t$, qui est bien une représentation \emph{AR d'ordre fini}
(et stationnaire car $0.66<1$).\end{reponse}
\begin{attention}
Sans simplifier, $\Theta(B)=1-B$ a une racine $B=1$ (non inversible) : on conclurait à tort
qu'il n'y a pas de représentation AR. Le facteur commun \emph{racine commune} aux deux
polynômes est la clef.
\end{attention}

\textbf{Q\,B.4 --- $X_t=1.6X_{t-1}-0.6X_{t-2}+a_t-0.6a_{t-1}$ stationnaire ?}
$\Phi(B)=1-1.6B+0.6B^2=(1-B)(1-0.6B)$ : racine $B=1$ (module $=1$, pas $>1$).
\begin{reponse}Réponse : \bonne{Faux}. La présence d'une racine unité ($B=1$) rend le
processus \emph{non stationnaire} (intégré).\end{reponse}
Le facteur $\Theta(B)=1-0.6B$ n'est pas commun avec $(1-B)$ : aucune simplification. On a en
fait un ARIMA$(1,1,1)$ : $(1-0.6B)\,\Delta X_t=(1-0.6B)a_t$\dots non, plus précisément
$(1-B)(1-0.6B)X_t=(1-0.6B)a_t$. Le $(1-0.6B)$ se simplifie ! $(1-B)X_t=a_t$ : c'est une
\textbf{marche aléatoire} $\Rightarrow$ non stationnaire. Conclusion inchangée : \bonne{Faux}.

\textbf{Q\,B.5 --- $(1-0.8B)X_t=(1+0.4B)a_t$.}
ARMA(1,1) : $\phi_1=0.8$ ($|0.8|<1$ : stationnaire/causal) ;
$\Theta(B)=1+0.4B$, racine $B=-1/0.4=-2.5$, $|{-2.5}|>1$ : inversible.
\begin{enumerate}[label=\alph*.]
\item \textbf{Représentation AR ?} Le MA étant inversible, on peut écrire
$a_t=\Theta(B)^{-1}\Phi(B)X_t$, soit un \bonne{$\mathrm{AR}(\infty)$ (oui, d'ordre infini)}.
\item \textbf{Représentation MA ?} L'AR étant stationnaire/causal,
$X_t=\Phi(B)^{-1}\Theta(B)a_t$, soit un \bonne{$\mathrm{MA}(\infty)$ (oui, d'ordre infini)}.
\end{enumerate}
(Les coefficients sont calculés en Q\,C.3.)

\textbf{Q\,B.6 --- $X_t=0.8X_{t-1}+a_t+0.8a_{t-1}$ stationnaire ?}
ARMA(1,1) avec $\Phi(B)=1-0.8B$ : racine $B=1.25>1$.
\begin{reponse}\bonne{Oui, stationnaire} (la partie AR détermine la stationnarité ; $|0.8|<1$).
\end{reponse}
La partie MA ($1+0.8B$, racine $-1.25$) est par ailleurs inversible : le processus est aussi
inversible.

\textbf{Q\,B.7 --- Nature de quatre processus.}
\begin{enumerate}[label=\arabic*),leftmargin=*]
\item $X_t=\eps_t-\eps_{t-3}$ : MA(3) (coefficients $0,0,-1$). $\E=0$ (\bonne{centré}),
finie $\Rightarrow$ \bonne{stationnaire} ; c'est un ARIMA$(0,0,3)$ \bonne{de type ARMA/MA}
(non intégré).
\item $X_t=\sigma^2-\eps_t^2$ : $\E(X_t)=\sigma^2-\E(\eps_t^2)=0$ (\bonne{centré}). Les
$\eps_t^2$ étant i.i.d., $X_t$ est une suite i.i.d.\ de variance constante et $\gamma(k)=0$
($k\ne0$) : \bonne{bruit blanc}, donc \bonne{stationnaire} (ARMA$(0,0)$).
\item $X_t-X_{t-2}=\eps_t$, i.e.\ $(1-B^2)X_t=\eps_t$, $(1-B^2)=(1-B)(1+B)$ : deux racines
unités $\Rightarrow$ \bonne{non stationnaire, intégré} (type ARIMA/SARIMA).
\item $X_t=2t+\eps_t-\eps_{t-2}$ : $\E(X_t)=2t$ dépend de $t$ $\Rightarrow$ \bonne{non centré,
non stationnaire}. C'est un processus \emph{TS} (tendance déterministe) : $X_t-2t$ est un
MA(2) stationnaire ; ce n'est pas un ARIMA pur (la non-stationnarité est déterministe, pas
stochastique).
\end{enumerate}

\textbf{Q\,B.8 --- $X_t-aX_{t-1}=\eps_t$, $0<|a|<1$.}
\begin{enumerate}[label=\arabic*),leftmargin=*]
\item C'est un \bonne{AR(1)}. Racine de $1-aB=0$ : $B=1/a$, $|1/a|>1$ car $|a|<1$ :
\bonne{causal et stationnaire}.
\item $X_t=(1-aB)^{-1}\eps_t=\sum_{j\ge0}a^j\eps_{t-j}$ (forme $\mathrm{MA}(\infty)$, causale).
\item Si $(\eps_t)$ i.i.d.\ gaussien, alors $(X_t)$, combinaison linéaire de gaussiennes, est
un \bonne{processus gaussien} (donc strictement stationnaire).
\end{enumerate}

\textbf{Q\,B.9 --- $X_t-aX_{t-1}=\eps_t-2b\eps_{t-1}+b^2\eps_{t-2}$.}
$\Theta(B)=1-2bB+b^2B^2=(1-bB)^2$.
\begin{enumerate}[label=\roman*),leftmargin=*]
\item ARMA(1,2). Causal/stationnaire $\iff |a|<1$ (racine $1/a$ de $\Phi$).
\item Inversible $\iff$ racine de $(1-bB)^2=0$, soit $B=1/b$, de module $>1$, i.e.\ $|b|<1$.
Dans ce cas $\eps_t=\Theta(B)^{-1}\Phi(B)X_t=(1-bB)^{-2}(1-aB)X_t
=\Big(\sum_{j\ge0}(j+1)b^jB^j\Big)(1-aB)X_t$.
\end{enumerate}

\textbf{Q\,B.10 --- $(1-B)(1-B^{12})$ stationnarise.}
\begin{enumerate}[label=\arabic*),leftmargin=*]
\item $Z_t=\alpha+\beta t+s_t+X_t$. L'opérateur $(1-B^{12})$ annule la saisonnalité de période
12 ($s_t=s_{t-12}\Rightarrow(1-B^{12})s_t=0$) et transforme $\alpha+\beta t$ en
$(1-B^{12})(\alpha+\beta t)=12\beta$ (constante). Puis $(1-B)$ annule cette constante :
$(1-B)(1-B^{12})(\alpha+\beta t)=0$ et $(1-B)(1-B^{12})s_t=0$. Il reste
$(1-B)(1-B^{12})X_t$, combinaison linéaire finie de $X$ stationnaire $\Rightarrow$
stationnaire (moyenne nulle, autocovariance indépendante de $t$).
\item $Z_t=(\alpha+\beta t)s_t+X_t$. Ici la saisonnalité \emph{module} une tendance :
$(\alpha+\beta t)s_t$ a une \og tendance saisonnière \fg{} d'ordre 1 multipliée par une
périodique. $(1-B^{12})\big[(\alpha+\beta t)s_t\big]=\big[(\alpha+\beta t)-(\alpha+\beta(t-12))\big]s_t
=12\beta s_t$ (car $s_t=s_{t-12}$), qui reste une tendance$\times$saison d'ordre 0. Une seconde
application $(1-B^{12})$ donne $0$. Donc $(1-B^{12})^2Z_t=(1-B^{12})^2X_t$, stationnaire.
\end{enumerate}

\textbf{Q\,B.11 --- Sous-échantillonnage d'un AR(1).}
$X_t=\alpha X_{t-1}+a_t$, $Y_t=X_{2t}$. Alors
$X_{2t}=\alpha X_{2t-1}+a_{2t}=\alpha(\alpha X_{2t-2}+a_{2t-1})+a_{2t}
=\alpha^2X_{2t-2}+\alpha a_{2t-1}+a_{2t}$, c'est-à-dire
\[
Y_t=\alpha^2 Y_{t-1}+b_t,\qquad b_t=\alpha a_{2t-1}+a_{2t}.
\]
$b_t$ est centré, $\Var(b_t)=(\alpha^2+1)\sigma^2$, et $\Cov(b_t,b_{t-1})=
\Cov(\alpha a_{2t-1}+a_{2t},\ \alpha a_{2t-3}+a_{2t-2})=0$ (indices tous distincts) : $b_t$ est
un \bonne{bruit blanc}. Donc $Y_t$ est un \bonne{AR(1) de paramètre $\phi=\alpha^2$}, avec
$|\alpha^2|<1$ : stationnaire et causal.

\textbf{Q\,B.12 --- $(1-\phi B^2)Y_t=\delta+(1+\theta B)a_t$.}
\begin{enumerate}[label=\arabic*),leftmargin=*]
\item $Y_t=(1-\phi B^2)^{-1}\delta+(1-\phi B^2)^{-1}(1+\theta B)a_t$. Comme
$(1-\phi B^2)^{-1}=\sum_{i\ge0}\phi^iB^{2i}$, on a $\frac{\delta}{1-\phi}$ pour la constante et,
pour la partie aléatoire, $\big(\sum_i\phi^iB^{2i}\big)(1+\theta B)a_t$. D'où les
$\psi_j$ : pour $j$ pair $j=2i$, $\psi_{2i}=\phi^i$ ; pour $j$ impair $j=2i+1$,
$\psi_{2i+1}=\theta\phi^i$. (Donc $\psi_0=1,\psi_1=\theta,\psi_2=\phi,\psi_3=\theta\phi,\dots$)
\item Stationnaire/causal $\iff$ racines de $1-\phi B^2=0$, $B=\pm1/\sqrt{\phi}$, de module
$>1$, soit \bonne{$|\phi|<1$} (avec $\phi>0$ pour des racines réelles, sinon complexes
conjuguées de module $1/\sqrt{|\phi|}$). Le terme MA $(1+\theta B)$ ne joue pas sur la
stationnarité ; l'inversibilité requiert $|\theta|<1$.
\item Avec $\E(Y)=\frac{\delta}{1-\phi}$ et la forme $\mathrm{MA}(\infty)$ :
$\gamma(0)=\sigma^2\sum_j\psi_j^2$ et $\gamma(1)=\sigma^2\sum_j\psi_j\psi_{j+1}$.
Or $\sum_j\psi_j\psi_{j+1}=\sum_i\psi_{2i}\psi_{2i+1}+\psi_{2i+1}\psi_{2i+2}
=\sum_i\phi^i\theta\phi^i+\theta\phi^i\phi^{i+1}=\theta(1+\phi)\sum_i\phi^{2i}
=\dfrac{\theta(1+\phi)}{1-\phi^2}=\dfrac{\theta}{1-\phi}$, et
$\sum_j\psi_j^2=\sum_i(\phi^{2i}+\theta^2\phi^{2i})=\dfrac{1+\theta^2}{1-\phi^2}$. D'où
\[
\rho(1)=\frac{\gamma(1)}{\gamma(0)}=\frac{\theta/(1-\phi)}{(1+\theta^2)/(1-\phi^2)}
=\frac{\theta(1+\phi)}{1+\theta^2}.
\]
\end{enumerate}

% =====================================================================
\section{Représentations \texorpdfstring{$\mathrm{MA}(\infty)$ et $\mathrm{AR}(\infty)$}{MA-infini et AR-infini}}
\markboth{Représentations MA-infini et AR-infini}{}

\begin{rappel}
Pour $\Phi(B)X_t=\Theta(B)a_t$ : la forme $\mathrm{MA}(\infty)$ $X_t=\Psi(B)a_t$ a des
coefficients $\psi_k$ déterminés par $\Phi(B)\Psi(B)=\Theta(B)$, et la forme
$\mathrm{AR}(\infty)$ $\Pi(B)X_t=a_t$ par $\Theta(B)\Pi(B)=\Phi(B)$ ($\psi_0=\pi_0=1$). On
identifie les coefficients de $B^k$ de proche en proche.
\end{rappel}

\textbf{Q\,C.1 --- AR(2) $X_t-0.9X_{t-1}+0.2X_{t-2}=a_t$, coefficients $\psi_k$.}
$\Phi(B)=1-0.9B+0.2B^2$, $\Theta(B)=1$. La relation $\Phi(B)\Psi(B)=1$ donne, pour $k\ge1$,
\[
\psi_k=0.9\,\psi_{k-1}-0.2\,\psi_{k-2}\quad(\psi_0=1,\ \psi_{-1}=0).
\]
\[
\psi_1=0.9,\quad \psi_2=0.9(0.9)-0.2(1)=0.61,\quad \psi_3=0.9(0.61)-0.2(0.9)=0.369 .
\]
\emph{Forme générale.} Racines de $\Phi$ : $0.2B^2-0.9B+1=0\Rightarrow B\in\{2,\ 2.5\}$,
d'où les inverses $0.5$ et $0.4$. Donc $\psi_k=A(0.5)^k+C(0.4)^k$ avec $\psi_0=A+C=1$,
$\psi_1=0.5A+0.4C=0.9$ : $A=5,\ C=-4$.
\begin{reponse}
$\psi_1=0.9,\ \psi_2=0.61,\ \psi_3=0.369$, et $\boxed{\psi_k=5\,(0.5)^k-4\,(0.4)^k}$.
\end{reponse}

\textbf{Q\,C.2 --- $X_t=0.6X_{t-1}+0.2X_{t-2}+a_t-0.5a_{t-1}$ : 4 premiers $\pi_k$.}
$\Phi(B)=1-0.6B-0.2B^2$, $\Theta(B)=1-0.5B$. La forme AR$(\infty)$ : $a_t=\Pi(B)X_t$ avec
$\Theta(B)\Pi(B)=\Phi(B)$. En identifiant $B^k$ ($\pi_0=1$) :
\[
\pi_k-0.5\,\pi_{k-1}=\phi_k^{\Phi},\quad\text{où }\phi_1^{\Phi}=-0.6,\ \phi_2^{\Phi}=0.2,\
\phi_k^{\Phi}=0\ (k\ge3).
\]
$\pi_1=-0.6+0.5(1)=-0.1$ ; $\pi_2=0.2+0.5(-0.1)... $ attention au signe :
$\pi_2-0.5\pi_1=0.2\Rightarrow\pi_2=0.2+0.5(-0.1)$. Or $\phi_2^{\Phi}=-(-0.2)$ ? Reprenons :
$\Phi(B)=1-0.6B-0.2B^2$ donc le coefficient de $B^2$ est $-0.2$. Ainsi
$\pi_2-0.5\pi_1=-0.2\Rightarrow\pi_2=-0.2+0.5(-0.1)=-0.25$ ;
$\pi_3-0.5\pi_2=0\Rightarrow\pi_3=0.5(-0.25)=-0.125$.
\begin{reponse}
$\pi_0=1,\ \pi_1=-0.1,\ \pi_2=-0.25,\ \pi_3=-0.125$.\\
Représentation : $a_t=X_t-0.1X_{t-1}-0.25X_{t-2}-0.125X_{t-3}-\dots$, soit
$X_t=0.1X_{t-1}+0.25X_{t-2}+0.125X_{t-3}+\dots+a_t$.
\end{reponse}

\textbf{Q\,C.3 --- $(1-0.8B)X_t=(1+0.4B)a_t$ : coefficients AR et MA.}
\emph{AR$(\infty)$} : $\Theta(B)\Pi(B)=\Phi(B)$ avec $\Theta=1+0.4B$, $\Phi=1-0.8B$ :
$\pi_k+0.4\pi_{k-1}=\phi_k^{\Phi}$ où $\phi_1^{\Phi}=-0.8$, $\phi_k^{\Phi}=0$ ($k\ge2$).
$\pi_1=-0.8-0.4(1)=-1.2$ ; $\pi_2=-0.4(-1.2)=0.48$ ; $\pi_3=-0.4(0.48)=-0.192$.
\begin{reponse}AR : $\pi_1=-1.2,\ \pi_2=0.48,\ \pi_3=-0.192$
(d'où $a_t=X_t-1.2X_{t-1}+0.48X_{t-2}-0.192X_{t-3}+\dots$).\end{reponse}
\emph{MA$(\infty)$} : $\Phi(B)\Psi(B)=\Theta(B)$ : $\psi_k-0.8\psi_{k-1}=\theta_k^{\Theta}$ où
$\theta_1^{\Theta}=0.4$, $\theta_k^{\Theta}=0$ ($k\ge2$).
$\psi_1=0.4+0.8=1.2$ ; $\psi_2=0.8(1.2)=0.96$ ; $\psi_3=0.8(0.96)=0.768$.
\begin{reponse}MA : $\psi_1=1.2,\ \psi_2=0.96,\ \psi_3=0.768$
(et $\psi_k=1.2\,(0.8)^{k-1}$ pour $k\ge1$).\end{reponse}

\textbf{Q\,C.4 --- $(1-\Phi_1B^7)x_t=a_t$ (AR saisonnier d'ordre 7).}
\begin{enumerate}[label=\alph*.]
\item $x_t=(1-\Phi_1B^7)^{-1}a_t=\sum_{j\ge0}\Phi_1^{\,j}a_{t-7j}$ : seuls les retards
multiples de 7 apparaissent ($\psi_{7j}=\Phi_1^{j}$, autres $\psi=0$).
\item $\gamma_0=\Var(x_t)=\sigma^2\sum_{j\ge0}\Phi_1^{2j}=\dfrac{\sigma^2}{1-\Phi_1^2}$
(stationnaire si $|\Phi_1|<1$).
\item Par Yule--Walker $\rho(k)=\Phi_1\rho(k-7)$, et $\rho(k)=0$ si $k$ n'est pas multiple
de 7. Donc $\rho_7=\Phi_1$, $\rho_{14}=\Phi_1^2$, et $\rho_1=0$, $\rho_8=0$ (8 non multiple
de 7).
\item La FAP d'un AR \og saisonnier pur \fg{} ne survit qu'aux retards $7,14,\dots$ jusqu'à
l'ordre du modèle (ici un seul terme) : $P_7=\Phi_1$, puis $P_{14}=0$ ; et $P_1=0$, $P_8=0$.
\end{enumerate}

\textbf{Q\,C.5 --- AR(2) $X_t-0.9X_{t-1}+0.2X_{t-2}=a_t$ (complet).}
\begin{enumerate}[label=\arabic*),leftmargin=*]
\item Racines $\{2,2.5\}$ de module $>1$ : \bonne{stationnaire et causal}.
\item Yule--Walker : $\rho_1=\dfrac{\phi_1}{1-\phi_2}=\dfrac{0.9}{1-(-0.2)}=\dfrac{0.9}{1.2}=0.75$ ;
$\rho_2=\phi_1\rho_1+\phi_2=0.9(0.75)-0.2=0.475$.
\item $\rho_k=0.9\rho_{k-1}-0.2\rho_{k-2}$ ($k\ge3$) ; forme générale
$\rho_k=A(0.5)^k+C(0.4)^k$ avec $\rho_0=1,\rho_1=0.75$ : $A+C=1$, $0.5A+0.4C=0.75\Rightarrow
A=3.5,\ C=-2.5$, soit $\rho_k=3.5(0.5)^k-2.5(0.4)^k$.
\item $\psi_k=5(0.5)^k-4(0.4)^k$ (cf.\ Q\,C.1).
\end{enumerate}

\textbf{Q\,C.6 --- Identifier $\phi,\theta$ depuis le début de la suite $\psi_k$.}
Pour $x_t=\phi x_{t-1}+a_t-\theta a_{t-1}$ (ARMA(1,1)), la $\mathrm{MA}(\infty)$ vérifie
$\psi_0=1$, $\psi_1=\phi-\theta$, et $\psi_k=\phi\,\psi_{k-1}$ pour $k\ge2$ (le rapport
$\psi_{k}/\psi_{k-1}=\phi$ devient constant).
\begin{recette}
\textbf{Méthode.} (i) Le rapport stabilisé $\psi_{k}/\psi_{k-1}$ pour $k\ge2$ donne
directement $\boxed{\phi}$. (ii) Puis $\theta=\phi-\psi_1$.
\end{recette}
Par exemple, si l'imprimante affiche $\psi_0=1,\ \psi_1=0.3,\ \psi_2=0.15,\ \psi_3=0.075,\dots$,
alors $\phi=\psi_2/\psi_1=0.5$ et $\theta=\phi-\psi_1=0.5-0.3=0.2$.
\emph{(Les valeurs numériques exactes dépendent de la suite imprimée sur le sujet ; appliquer
la méthode encadrée.)}

% =====================================================================
\section{Autocorrélations simples (ACF) et partielles (PACF)}
\markboth{Autocorrélations ACF / PACF}{}

\textbf{Q\,D.1 --- AR(2), $(1-0.85B+0.32B^2)$, $\rho(1)=0.84,\rho(2)=0.69,\rho(3)=0.52$.}
On lit $\phi_1=0.85$, $\phi_2=-0.32$ (car $\Phi(B)=1-\phi_1B-\phi_2B^2=1-0.85B+0.32B^2$).
\begin{enumerate}[label=\arabic*),leftmargin=*]
\item Pour un AR(2), la FAP s'annule au-delà de 2 et $P_2=\phi_2$. De plus $P_1=\rho_1$.
\begin{reponse}$\phi_{11}=P_1=\rho(1)=0.84$ \quad et\quad $\phi_{22}=P_2=\phi_2=-0.32$.\end{reponse}
\item Yule--Walker : $\rho(1)=\dfrac{\phi_1}{1-\phi_2}$ et $\rho(2)=\phi_1\rho(1)+\phi_2$.
\item Numériquement : $\rho(1)=\dfrac{0.85}{1.32}=0.6439$ et
$\rho(2)=0.85(0.6439)-0.32=0.2273$.
\end{enumerate}
\begin{remarque}
Les valeurs \emph{empiriques} annoncées ($0.84,0.69$) ne coïncident pas avec les valeurs
\emph{théoriques} ($0.644,0.227$) issues de l'AR(2) estimé : c'est le sens de la question
(incohérence entre l'ACF observée et le modèle ajusté).
\end{remarque}

\textbf{Q\,D.2 --- $(1-0.8B^2)X_t=a_t$.}
AR(2) avec $\phi_1=0,\ \phi_2=0.8$. Yule--Walker :
$\rho_1=\dfrac{\phi_1}{1-\phi_2}=0$ ; $\rho_2=\phi_1\rho_1+\phi_2=0.8$.
FAP : $P_1=\rho_1=0$ ; $P_2=\dfrac{\rho_2-\rho_1^2}{1-\rho_1^2}=0.8=\phi_2$.
\begin{reponse}$\rho_1=0,\ \rho_2=0.8$ ; $P_1=0,\ P_2=0.8$.\end{reponse}

\textbf{Q\,D.3 --- $X_t=0.6X_{t-1}+0.8X_{t-2}-0.3X_{t-3}+a_t$, 3 premières $\rho_k$.}
AR(3) : $\phi_1=0.6,\phi_2=0.8,\phi_3=-0.3$. Yule--Walker ($\rho_0=1$, $\rho_{-k}=\rho_k$) :
\begin{align*}
\rho_1&=\phi_1+\phi_2\rho_1+\phi_3\rho_2,\\
\rho_2&=\phi_1\rho_1+\phi_2+\phi_3\rho_1,\\
\rho_3&=\phi_1\rho_2+\phi_2\rho_1+\phi_3 .
\end{align*}
Les deux premières donnent le système
$\begin{cases}(1-\phi_2)\rho_1-\phi_3\rho_2=\phi_1\\(\phi_1+\phi_3)\rho_1-\rho_2=-\phi_2\end{cases}$,
soit $\begin{cases}0.2\rho_1+0.3\rho_2=0.6\\0.3\rho_1-\rho_2=-0.8\end{cases}$.
On résout : $\rho_1\approx1.1304$ ?? \emph{Attention} : ce processus n'est pas stationnaire
(voir ci-dessous), donc les équations de Yule--Walker n'ont pas nécessairement de solution
admissible $|\rho|\le1$.
\begin{attention}
Vérifions d'abord la stationnarité : $\Phi(B)=1-0.6B-0.8B^2+0.3B^3$. En $B=1$ :
$1-0.6-0.8+0.3=-0.1\ne0$. Les racines doivent être de module $>1$. Une étude numérique montre
qu'une racine est de module $<1$ : \textbf{le processus n'est pas stationnaire}, l'ACF
théorique n'est donc pas définie au sens usuel. Si l'énoncé suppose la stationnarité, on
applique néanmoins formellement Yule--Walker ; sinon, on conclut à la non-stationnarité.
\end{attention}
\emph{Application formelle de Yule--Walker} (si l'on suppose la récurrence valide) :
$\rho_k=0.6\rho_{k-1}+0.8\rho_{k-2}-0.3\rho_{k-3}$, avec $\rho_0=1$ et la résolution du système
$2\times2$ ci-dessus pour $(\rho_1,\rho_2)$, puis $\rho_3=0.6\rho_2+0.8\rho_1-0.3$.

\textbf{Q\,D.4 --- AR(2) $X_t-0.9X_{t-1}+0.2X_{t-2}=a_t$ : $\rho_1,\rho_2,\rho_k$.}
(Identique à Q\,C.5) : $\rho_1=0.75$, $\rho_2=0.475$, et pour $k\ge3$
$\rho_k=0.9\rho_{k-1}-0.2\rho_{k-2}=3.5(0.5)^k-2.5(0.4)^k$.

\textbf{Q\,D.5 --- Identification ARMA à partir d'un tableau ACF/PACF.}
\begin{center}\small
\begin{tabular}{cccccc}
\toprule $k$&1&2&3&4&5\\\midrule
$\rho_k$&0.65&0.47&0.35&0.24&0.18\\
$P_k$&0.65&$-0.31$&$-0.08$&$-0.15$&0.10\\\bottomrule
\end{tabular}
\end{center}
L'ACF \emph{décroît} régulièrement (pas de coupure nette) ; la PACF est \emph{forte en 1}
($0.65$) puis devient faible/négligeable dès le rang 2 (les $|P_k|\lesssim2/\sqrt{100}=0.2$
sont à la limite de significativité).
\begin{reponse}On retient un \bonne{AR(1)} : ACF qui décroît géométriquement
($\rho_k\approx0.65^{\,k}$ : $0.65,0.42,0.27,\dots$, proche des valeurs observées) et PACF qui
\og coupe \fg{} après le rang 1.\end{reponse}

\textbf{Q\,D.6 --- $z_t=x_t+y_t$, $x_t\sim$AR(1), $y_t\sim$MA(1) indépendants.}
La somme de deux processus ARMA indépendants ARMA$(p_1,q_1)$ et ARMA$(p_2,q_2)$ est un
ARMA$(p_1+p_2,\ \max(p_1+q_2,\ p_2+q_1))$. Ici $x_t$ est ARMA(1,0), $y_t$ ARMA(0,1) :
$p=1+0=1$, $q=\max(1+1,\ 0+0)=2$.
\begin{reponse}$z_t$ est un \bonne{ARMA$(1,1)$} en général (l'ordre MA \emph{apparent} est $2$,
mais il se réduit) --- plus précisément, $(1-\phi B)z_t=(1-\phi B)x_t+(1-\phi B)y_t
=a_t+(1-\phi B)(1+\theta B)b_t$, dont la partie MA est d'ordre $\le2$ : $z_t$ est un
\bonne{ARMA$(1,2)$} qui, après calcul des autocovariances, est équivalent à un
\bonne{ARMA$(1,1)$}.\end{reponse}
Réponse attendue au QCM : \bonne{ARMA$(p,q)$} (ici $p=1$, $q\le2$).

% =====================================================================
\section{Prévision et intervalles de confiance}
\markboth{Prévision et intervalles de confiance}{}

\begin{rappel}
Prévision optimale (espérance conditionnelle) : on remplace les chocs \emph{futurs} par $0$,
les chocs \emph{passés} par les résidus, et les $X$ futurs par leurs prévisions.
\emph{Erreur de prévision} à horizon $k$ : $\Delta_k=X_{n+k}-\Xhat_{n+k}=\sum_{j=0}^{k-1}
\psi_j a_{n+k-j}$, $\Var(\Delta_k)=\sigma^2\sum_{j=0}^{k-1}\psi_j^2$.
IC à $95\%$ : $\Xhat_{n+k}\pm1.96\,\sigma\sqrt{\sum_{j=0}^{k-1}\psi_j^2}$.
\end{rappel}

\textbf{Q\,E.1 --- MA(1) $X_t=2+a_t+0.8a_{t-1}$, $\sigma^2=2$, $a_t=1.6$, $a_{t-1}=-3.4$.}
$\Xhat_{t+1}=2+0.8\,a_t=2+0.8(1.6)=\mathbf{3.28}$.
$\Xhat_{t+2}=2$ (horizon $\ge2$ pour un MA(1) : prévision $=\mu=2$).
$\Var(\Delta_1)=\sigma^2=2$ ; $\Var(\Delta_2)=(1+\theta^2)\sigma^2=(1+0.64)\cdot2=3.28$.
\begin{reponse}
$\Xhat_{t+1}=3.28$, IC$_{95\%}=3.28\pm1.96\sqrt2=[0.51,\ 6.05]$.\\
$\Xhat_{t+2}=2$, IC$_{95\%}=2\pm1.96\sqrt{3.28}=[-1.55,\ 5.55]$.
\end{reponse}

\textbf{Q\,E.2 --- $(1-0.8B)X_t=2+a_t$, $a_T=1.6$, $X_{T-1}=3.0$, $\sigma_a^2=1$.}
Modèle $X_t=2+0.8X_{t-1}+a_t$. D'abord $X_T=2+0.8(3.0)+1.6=6.0$.
$\Xhat_{T+1}=2+0.8X_T=2+4.8=\mathbf{6.8}$ ; $\Xhat_{T+2}=2+0.8\Xhat_{T+1}=\mathbf{7.44}$.
$\Var(\Delta_1)=\sigma^2=1$ ; $\Var(\Delta_2)=(1+\phi^2)\sigma^2=1.64$.
\begin{reponse}
$\Xhat_{T+1}=6.8$, IC $=6.8\pm1.96=[4.84,\ 8.76]$.\\
$\Xhat_{T+2}=7.44$, IC $=7.44\pm1.96\sqrt{1.64}=[4.93,\ 9.95]$.
\end{reponse}

\textbf{Q\,E.3 --- MA(2) $Y_t=a_t+0.8a_{t-1}+0.4a_{t-2}$, résidus connus.}
$\Xhat_{101}=0.8a_{100}+0.4a_{99}=0.8(0.6)+0.4(-0.3)=\mathbf{0.36}$.
$\Xhat_{102}=0.4a_{100}=0.4(0.6)=\mathbf{0.24}$.
$\Xhat_{103}=0$ (horizon $3>q=2$ : prévision $=\mu=0$).
$\Var(\Delta_1)=\sigma^2=1$ ; $\Var(\Delta_2)=(1+0.8^2)=1.64$ ;
$\Var(\Delta_3)=(1+0.8^2+0.4^2)=1.8$.
\begin{reponse}
IC : $0.36\pm1.96=[-1.60,2.32]$ ; $0.24\pm1.96\sqrt{1.64}=[-2.27,2.75]$ ;
$0\pm1.96\sqrt{1.8}=[-2.63,2.63]$.
\end{reponse}

\textbf{Q\,E.4 --- MA(1) (sortie EViews), série 06--09/2012.}
Modèle $X_t=C+a_t+\theta a_{t-1}$ avec $C=1.99678$, $\theta=-0.492396$, $\sigma^2=1$.
On reconstitue les résidus par $a_t=X_t-C-\theta a_{t-1}$ (en posant $a_{05/2012}=0$) :
\[
a_{06}=-0.7968,\quad a_{07}=1.1109,\quad a_{08}=2.7503,\quad a_{09}=5.7574 .
\]
$\Xhat_{10}=C+\theta a_{09}=1.99678-0.492396(5.7574)=\mathbf{-0.838}$ ;
$\Xhat_{11}=\Xhat_{12}=C=1.99678$ (horizon $\ge2$).
$\Var(\Delta_1)=1$ ; $\Var(\Delta_{k\ge2})=(1+\theta^2)=1.2425$.
\begin{reponse}
$\Xhat_{10}=-0.838$, IC $=[-2.80,\,1.12]$ ; $\Xhat_{11}=\Xhat_{12}=1.997$,
IC $=1.997\pm1.96\sqrt{1.2425}=[-0.19,\,4.18]$.
\end{reponse}
\begin{attention}
La série $1.2;3.5;4.2;6.4$ est nettement \emph{croissante} (tendance) : un MA(1) y est
inadapté (d'où une prévision aberrante). L'exercice illustre l'importance de stationnariser
avant de modéliser ; on présente ici l'application mécanique de la formule.
\end{attention}

\textbf{Q\,E.5 --- ARMA(1,1) (Série 2, Ex.~1).}
$X_t=0.495639\,X_{t-1}+a_t+0.413047\,a_{t-1}$ (sans constante). Résidus reconstitués
($a_{2007}=0$) : $a_{08}=1.1565$, $a_{09}=1.7806$, $a_{10}=-0.1215$, $a_{11}=1.3598$.
$\Xhat_{2012}=0.495639(2.4)+0.413047(1.3598)=\mathbf{1.751}$ ;
$\Xhat_{2013}=0.495639\,\Xhat_{2012}=\mathbf{0.868}$ (au-delà de l'horizon 1, $a$ futur $=0$).
\begin{reponse}$\Xhat_{2012}\approx1.75$, $\Xhat_{2013}\approx0.87$.\end{reponse}

\textbf{Q\,E.6 --- AR(1) à moyenne (Série 2, Ex.~2).}
$C=\mu=7.124608$, $\phi=0.297718$, $\sigma^2=1$, $X_{2011}=10.4$. Forme centrée
$X_t-\mu=\phi(X_{t-1}-\mu)+a_t$ :
\[
\Xhat_{2012}=\mu+\phi(X_{2011}-\mu)=\mathbf{8.10},\quad
\Xhat_{2013}=\mu+\phi(\Xhat_{2012}-\mu)=\mathbf{7.42},\quad
\Xhat_{2014}=\mathbf{7.21}.
\]
$\Var(\Delta_k)=\sigma^2\sum_{j=0}^{k-1}\phi^{2j}$ : $1$ ; $1.0886$ ; $1.0965$.
\begin{reponse}
IC : $8.10\pm1.96=[6.14,10.06]$ ; $7.42\pm1.96\sqrt{1.0886}=[5.37,9.46]$ ;
$7.21\pm1.96\sqrt{1.0965}=[5.16,9.26]$.
\end{reponse}

\textbf{Q\,E.7 --- AR(1) (brouillon), constante $5.0552$, $\phi=0.29$.}
On interprète le modèle comme $X_t=5.0552+0.29X_{t-1}+a_t$ (intercept), $\sigma^2=1$ par
défaut, dernière observation $X_{2011}=7.12$.
$\Xhat_{2012}=5.0552+0.29(7.12)=7.120$ ; $\Xhat_{2013}=5.0552+0.29(7.120)=7.120$ ;
$\Xhat_{2014}\approx7.120$ (le processus converge vite vers sa moyenne
$\mu=5.0552/(1-0.29)=7.12$).
$\Var(\Delta_k)=\sum_{j=0}^{k-1}0.29^{2j}$ : $1,\ 1.084,\ 1.091$.
\begin{reponse}
$\Xhat\approx7.12$ pour les trois horizons ; IC $\approx7.12\pm1.96\sqrt{\Var(\Delta_k)}$.
\emph{(Énoncé reconstitué : valeurs sous réserve des hypothèses ci-dessus.)}
\end{reponse}

% =====================================================================
\section{Intégration, ARIMA et tests de Dickey--Fuller}
\markboth{Intégration, ARIMA, Dickey--Fuller}{}

\textbf{Q\,F.1 --- $X_t=2.65X_{t-1}-2.3X_{t-2}+0.65X_{t-3}+a_t$.}
$\Phi(B)=1-2.65B+2.3B^2-0.65B^3$. En $B=1$ : $1-2.65+2.3-0.65=0$ : racine unité. On
factorise par $(1-B)$ :
\[
\Phi(B)=(1-B)(1-1.65B+0.65B^2)=(1-B)(1-B)(1-0.65B)=(1-B)^2(1-0.65B).
\]
(Vérification du second facteur en $B=1$ : $1-1.65+0.65=0$, d'où la deuxième racine unité.)
\begin{enumerate}[label=\alph*.]
\item \bonne{$X_t$ est intégré d'ordre $d=2$} ($\sim I(2)$).
\item $(1-B)^2(1-0.65B)X_t=a_t$ ; en posant $W_t=(1-B)^2X_t$ on a $(1-0.65B)W_t=a_t$ (AR(1)
stationnaire). L'écriture Dickey--Fuller pour $y_t=\Delta x_t$ fait apparaître, après une
différenciation, une racine unité restante (le processus $\Delta x_t$ est encore $I(1)$).
\item \begin{reponse}Réponse : \bonne{ARIMA$(1,2,0)$}.\end{reponse}
(partie AR d'ordre 1 : $1-0.65B$ ; $d=2$ ; pas de partie MA).
\end{enumerate}

\textbf{Q\,F.2 --- $x_t=2.7x_{t-1}-2.4x_{t-2}+0.7x_{t-3}+a_t$.}
$\Phi(B)=1-2.7B+2.4B^2-0.7B^3$. En $B=1$ : $1-2.7+2.4-0.7=0$. Factorisation :
$\Phi(B)=(1-B)(1-1.7B+0.7B^2)=(1-B)^2(1-0.7B)$ (le facteur $1-1.7B+0.7B^2$ s'annule aussi en
$B=1$).
\begin{reponse}\bonne{$x_t\sim I(2)$}, gouverné par un \bonne{ARIMA$(1,2,0)$}
($\Delta^2x_t=0.7\,\Delta^2x_{t-1}+a_t$).\end{reponse}

\textbf{Q\,F.3 --- $X_t=1.6X_{t-1}-0.3X_{t-2}+0.6X_{t-3}+a_t$ intégré ?}
$\Phi(B)=1-1.6B+0.3B^2-0.6B^3$. En $B=1$ : $1-1.6+0.3-0.6=-0.9\ne0$ : pas de racine unité.
\begin{reponse}Réponse : \bonne{Non intégré} (c'est un AR(3) ; sa stationnarité dépend du
module des racines, mais il n'y a pas de racine $=1$).\end{reponse}

\textbf{Q\,F.4 --- Série E : MCO et test de Dickey--Fuller.}
\begin{enumerate}[label=\alph*.]
\item MCO sans constante : $\phihat=\dfrac{\sum Y_tY_{t-1}}{\sum Y_{t-1}^2}
=\dfrac{25\,842.3}{25\,943.8}=\mathbf{0.99609}$.
\item $\widehat{\sigma}_{\phihat}=\sqrt{\dfrac{\hat\sigma_a^2}{\sum Y_{t-1}^2}}
=\sqrt{\dfrac{0.954}{25\,943.8}}=\mathbf{0.006064}$.
\item Statistique DF : $t_\phi=\dfrac{\phihat-1}{\widehat{\sigma}_{\phihat}}
=\dfrac{0.99609-1}{0.006064}=\mathbf{-0.645}$. Comme $t_\phi=-0.645$ est
\emph{supérieur} à la valeur critique $-2.89$ (donc \emph{hors} de la région de rejet), on
\bonne{ne rejette pas $H_0:\phi=1$} : la série possède une racine unité, elle est
\bonne{non stationnaire ($I(1)$)}.
\item Après différenciation, $\rhohat(1)=-0.035\approx0$ : $\Delta Y_t$ est un \emph{bruit
blanc}. Donc $Y_t$ est une \bonne{marche aléatoire}, modèle \bonne{ARIMA$(0,1,0)$}.
\item $s_{\Delta Y}^2=0.954 < s_Y^2=86.747$ : la différenciation a \emph{fortement réduit} la
variance (passage d'une série $I(1)$ à un bruit blanc). Mais la différenciation ne réduit
\emph{pas toujours} la variance : sur une série \emph{déjà stationnaire}, sur-différencier
\emph{augmente} la variance (introduction d'une racine MA unité). On ne différencie que si un
test de racine unité le justifie.
\end{enumerate}

\textbf{Q\,F.5 --- Choix du générateur (300 obs.).}
\begin{itemize}[leftmargin=*]
\item (a) $Y_t=1.8Y_{t-1}+a_t$ : $|\phi|=1.8>1$ $\Rightarrow$ \emph{explosif} (croissance
non bornée). À écarter si la trajectoire reste bornée.
\item (b) $Y_t=0.8Y_{t-1}+10+a_t$ : AR(1) stationnaire de moyenne $\mu=10/(1-0.8)=50$ et de
variance $\sigma^2/(1-0.64)=2.78$ ($\sigma=1$). Série centrée \emph{autour de 50}.
\item (c) $Y_t=0.8Y_{t-1}+0.5a_t$ : AR(1) stationnaire, moyenne $0$, écart-type
$\sqrt{0.25/0.36}=0.83$ : faibles amplitudes autour de 0.
\item (d) $Y_t=0.8Y_{t-1}+5a_t$ : AR(1) stationnaire, moyenne $0$, écart-type
$\sqrt{25/0.36}=8.33$ : fortes amplitudes autour de 0.
\end{itemize}
\begin{reponse}
On choisit le modèle dont la \emph{moyenne} et l'\emph{amplitude} correspondent au graphique :
si la série oscille autour de $50$ $\to$ (b) ; autour de $0$ avec faible amplitude $\to$ (c) ;
autour de $0$ avec forte amplitude $\to$ (d) ; si elle explose $\to$ (a). On écarte (a) dès que
la série est bornée, et on tranche entre (b),(c),(d) par le niveau moyen et l'écart-type
observés.
\end{reponse}

% =====================================================================
\section{Estimation et équations de Yule--Walker}
\markboth{Estimation et Yule--Walker}{}

\textbf{Q\,G.1 --- $x_t=0.4x_{t-1}-0.2x_{t-2}+a_t$, $\sigma_a^2=12.8$.}
\begin{enumerate}[label=\arabic*),leftmargin=*]
\item $\Phi(B)=1-0.4B+0.2B^2$ ; discriminant $0.16-0.8<0$, racines complexes de module
$1/\sqrt{0.2}=2.236>1$ : \bonne{stationnaire}.
\item $\E(x_t)=\dfrac{\mu}{1-\phi_1-\phi_2}=0$ (pas de constante).
\item Yule--Walker : $\rho_1=\dfrac{\phi_1}{1-\phi_2}=\dfrac{0.4}{1.2}=0.3333$ ;
$\rho_2=\phi_1\rho_1+\phi_2=0.4(0.3333)-0.2=-0.0667$ ;
$\rho_3=\phi_1\rho_2+\phi_2\rho_1=0.4(-0.0667)-0.2(0.3333)=-0.0933$.\\
Variance : $\gamma_0=\dfrac{\sigma_a^2}{1-\phi_1\rho_1-\phi_2\rho_2}
=\dfrac{12.8}{1-0.4(0.3333)+0.2(-0.0667)}=\dfrac{12.8}{0.8533}=\mathbf{15.0}$.\\
FAP : $P_1=\rho_1=0.3333$ ; $P_2=\dfrac{\rho_2-\rho_1^2}{1-\rho_1^2}=\phi_2=-0.2$ ; $P_3=0$
(AR(2) $\Rightarrow$ FAP nulle au-delà de 2).
\end{enumerate}

\textbf{Q\,G.2 --- AR(1) $Y_t=0.8Y_{t-1}+\eps_t$, $\sigma_\eps^2=2$.}
\begin{enumerate}[label=\arabic*),leftmargin=*]
\item $|\phi|=0.8<1$ : \bonne{stationnaire}. Un AR est \emph{toujours inversible} : \bonne{oui}.
\item $\gamma_0=\dfrac{\sigma^2}{1-\phi^2}=\dfrac{2}{0.36}=5.5556$ ; $\gamma_k=\phi^k\gamma_0$ :
$\gamma_1=4.444,\ \gamma_2=3.556,\ \gamma_3=2.844,\ \gamma_4=2.276,\ \gamma_5=1.820$.
\item $\rho_k=\phi^k$ : $0.8,\ 0.64,\ 0.512,\ 0.4096,\ 0.3277$.
\item $Y_t=\sum_{j\ge0}\phi^j\eps_{t-j}$ donc $\theta_j=\phi^j=0.8^j$ :
$\theta_1=0.8,\theta_2=0.64,\theta_3=0.512,\theta_4=0.4096,\theta_5=0.32768$.
\end{enumerate}

\textbf{Q\,G.3 --- Estimation de $a$ par Yule--Walker, AR(1).}
\begin{enumerate}[label=\alph*.]
\item $\gammahat(k)=\dfrac1n\sum_{j=1}^{n-k}(x_j-\bar x)(x_{j+k}-\bar x)$.
\item Yule--Walker AR(1) : $\rho(1)=a$, soit $\gamma(1)=a\,\gamma(0)$.
\item $\hat a=\rhohat(1)=\dfrac{\gammahat(1)}{\gammahat(0)}$.
\item Sous $(\eps_t)$ i.i.d.$(0,\sigma^2)$, $\hat a$ est \emph{convergent} et asymptotiquement
normal : $\sqrt n(\hat a-a)\xrightarrow{d}N\big(0,\,1-a^2\big)$.
\item Pour $a=1/2$ : $\Xhat_{t+1}=\tfrac12X_t$, $\Xhat_{t+2}=\tfrac14X_t$, \dots,
$\Xhat_{t+h}=(\tfrac12)^hX_t$ (prévision optimale d'un AR(1) causal, qui décroît vers la
moyenne $0$).
\end{enumerate}

\textbf{Q\,G.4 --- Identification à partir d'une sortie EViews.}
On lit l'ACF et la PACF empiriques (jusqu'au retard 10) : si l'ACF décroît et la PACF coupe à
l'ordre $p$ $\to$ AR$(p)$ ; si l'ACF coupe à $q$ et la PACF décroît $\to$ MA$(q)$ ; si les deux
décroissent $\to$ ARMA. On estime ensuite par Yule--Walker (AR) ou maximum de vraisemblance,
et on \emph{valide} par Ljung--Box sur les résidus. \emph{(Réponse dépendant des graphiques du
sujet ; appliquer la grille de lecture du chapitre Box--Jenkins.)}

% =====================================================================
\section{Décomposition : tendance et saisonnalité}
\markboth{Décomposition : tendance, saisonnalité}{}

\textbf{Q\,H.1 --- Cours d'une action sur 5 jours.}
Régression linéaire $y=at+b$ avec $t=1,\dots,5$, $y=(135,143,140,154,152)$.
$\bar t=3,\ \bar y=144.8$ ; $\hat a=\dfrac{\sum(t-\bar t)(y-\bar y)}{\sum(t-\bar t)^2}
=\dfrac{45}{10}=\mathbf{4.5}$ ; $\hat b=\bar y-\hat a\bar t=144.8-13.5=\mathbf{131.3}$.
\begin{reponse}Tendance : $\hat y_t=4.5\,t+131.3$. Prévision $t=6$ :
$\hat y_6=4.5(6)+131.3=\mathbf{158.3}$.\end{reponse}

\textbf{Q\,H.2 --- Ventes trimestrielles (Série 1, Ex.~1).}
\emph{Méthode (modèle additif, tendance par MCO).}
\begin{enumerate}[label=\alph*),leftmargin=*]
\item Tracé : tendance croissante + saisonnalité (pic au T4, creux au T3).
\item Tendance : régression $Z_t=at+b$ sur $t=1,\dots,16$. Avec $\sum t=136$, $\sum Z=4731$,
$\sum t^2=1496$, $\sum tZ=44\,330$ : $\hat a=\dfrac{16(44330)-136(4731)}{16(1496)-136^2}
=\dfrac{709280-643416}{23936-18496}=\dfrac{65864}{5440}=12.11$ ;
$\hat b=\dfrac{4731-12.11(136)}{16}=193.81$. Donc $\hat Z_t=12.11\,t+193.81$.
\item Coefficients saisonniers : calculer $e_t=Z_t-\hat Z_t$, moyenner par trimestre, puis
recentrer (somme nulle en additif). \item Prévisions : $\hat Z_{t}=12.11t+193.81+S_{j(t)}$
pour $t=17,\dots,20$. \item CVS : $Z_t-S_{j(t)}$.
\item Interprétation : $S_j$ mesure l'écart systématique du trimestre $j$ à la tendance.
\end{enumerate}
\emph{(Le détail numérique complet des 16 valeurs figure dans le cours « décomposition ». La
démarche est la pièce évaluée.)}

\textbf{Q\,H.3 --- Lettres reçues (Série 1, Ex.~2), MM d'ordre 4.}
\emph{Modèle multiplicatif} (amplitude croissant avec le niveau).
\begin{enumerate}[label=\alph*),leftmargin=*]
\item Saisonnalité marquée (T2 fort, T4 faible) avec légère tendance.
\item Tendance par MM4 \emph{centrée} : $T_t=\frac18 X_{t-2}+\frac14(X_{t-1}+X_t+X_{t+1})
+\frac18X_{t+2}$ (deux valeurs perdues à chaque bord).
\item Coefficients saisonniers : rapports $X_t/T_t$ moyennés par trimestre puis normalisés
(somme $=4$).
\item CVS : $X_t/S_{j(t)}$ ; on commente la disparition des pics saisonniers.
\end{enumerate}

\textbf{Q\,H.4 --- Indice de production industrielle 1976--1982 (MCO).}
\begin{enumerate}[label=\arabic*),leftmargin=*]
\item $\hat a=\dfrac{n\sum tY-\sum t\sum Y}{n\sum t^2-(\sum t)^2}
=\dfrac{28(53500)-406(3647)}{28(7714)-406^2}=\dfrac{1498000-1480682}{215992-164836}
=\dfrac{17318}{51156}=\mathbf{0.3385}$.\\
$\hat b=\dfrac{\sum Y-\hat a\sum t}{n}=\dfrac{3647-0.3385(406)}{28}=\mathbf{125.34}$.
\item Coefficients saisonniers : calculer $\tilde\eps_t=Y_t-\hat Y_t$ ($\hat Y_t=\hat at+\hat b$),
moyenner par trimestre, recentrer (additif, somme nulle).
\item Le tableau se complète colonne par colonne : $t$, $Y_t$, $\hat Y_t=0.3385t+125.34$,
$\tilde\eps_t=Y_t-\hat Y_t$, trimestre.
\end{enumerate}

% =====================================================================
\section{Identification graphique}
\markboth{Identification graphique}{}

\textbf{Q\,I.1 --- Associer chaque série à son graphique.}
\begin{itemize}[leftmargin=*]
\item $y_t=\eps_t$ : bruit pur, oscillations centrées sans structure.
\item $y_t=s_t+\eps_t$ : oscillations \emph{périodiques} autour de $0$ (saisonnalité additive,
moyenne nulle).
\item $y_t=a+s_t+\eps_t$ : saisonnalité périodique autour d'un \emph{niveau $a>0$}.
\item $y_t=a\,s_t\,\eps_t$ : amplitude \emph{modulée} (multiplicatif), variabilité croissante
avec le niveau.
\item $y_t=f_t+\eps_t$ : ondulation lente déterministe (sinusoïde) $+$ bruit additif.
\item $y_t=f_t\,\eps_t$ : sinusoïde dont l'amplitude du bruit suit $f_t$ (effet multiplicatif).
\end{itemize}
On distingue \emph{additif} (amplitude constante) de \emph{multiplicatif} (amplitude
proportionnelle au niveau), et la présence/absence d'un niveau moyen $a$ et d'une tendance
lente $f_t$.

\textbf{Q\,I.2 --- Identifier les processus derrière les ACF (Séries D, E, F).}
On utilise la grille du cours :
\begin{itemize}[leftmargin=*]
\item ACF qui \emph{décroît géométriquement} (sans coupure) $\to$ AR (le signe alterné
indique $\phi<0$).
\item ACF qui \emph{s'annule brutalement} après le retard $q$ $\to$ MA$(q)$.
\item ACF qui \emph{décroît très lentement} ($\approx1$) $\to$ processus \emph{non
stationnaire / intégré}.
\end{itemize}
\emph{(L'attribution précise D/E/F dépend des trois graphiques ; appliquer cette grille.)}

\textbf{Q\,I.3 --- Identification ACF/PACF (examen 17/12/2019).}
\begin{enumerate}[label=\arabic*),leftmargin=*]
\item PACF coupe à 1, ACF décroît $\to$ \bonne{AR(1)} ; si ACF coupe à 1 et PACF décroît
$\to$ MA(1) ; deux coupures à 2 $\to$ MA(2)/AR(2) selon laquelle coupe.
\item ACF qui décroît \emph{très lentement} $\to$ \bonne{ARIMA$(1,1,0)$} (racine unité).
\item PACF coupe à 2 $\to$ \bonne{AR(2)} ; coupe à 3 $\to$ AR(3) (selon le dernier pic
significatif).
\item Pics aux retards saisonniers $\to$ partie SARIMA ; choisir l'ordre saisonnier d'après
les pics aux multiples de la période $s=10$.
\item Pic unique au retard saisonnier $\to$ \bonne{SARIMA$(1,0,0)\times(1,0,0)_{12}$} (un
terme AR régulier $+$ un terme AR saisonnier).
\end{enumerate}
\emph{(Réponses indicatives : la justification s'appuie sur la coupure de l'ACF/PACF et la
position des pics saisonniers.)}

% =====================================================================
\section{Exercices de synthèse}
\markboth{Exercices de synthèse}{}

\textbf{Q\,J.1 --- AR(1) $X_t=\phi X_{t-1}+\eps_t$.}
\begin{enumerate}[label=\arabic*),leftmargin=*]
\item Affirmation \emph{fausse} : \bonne{(D)} \og $X_t$ et $\eps_t$ corrélés \fg{} est
\emph{l'affirmation correcte} ; la question demande celle qui n'est \emph{pas} correcte. Or
$\Cov(X_t,\eps_t)=\sigma^2\ne0$ (donc D \emph{est} vraie). Les non corrélations
$\Cov(\eps_t,X_{t-1})=0$ (A) et $\Cov(\eps_t,\eps_{t-1})=0$ (C) sont vraies ; en revanche
\bonne{(B) $X_t$ et $\eps_{t-1}$ \emph{sont} corrélés} (car $X_t$ contient $\eps_{t-1}$ via
$\phi\eps_{t-1}$). L'affirmation \emph{incorrecte} est donc \bonne{(B)}.
\item Correcte : \bonne{(A) $\cor(X_t,X_{t-1})=\phi$} (pour un AR(1), $\rho(1)=\phi$). (C) est
fausse car $\cor(X_t,\eps_t)\ne0$.
\item \bonne{(C) $\Var(X_t)=\dfrac{1}{1-\phi^2}$} (avec $\sigma_\eps^2=1$).
\item \bonne{(B)} : $Y_t=X_t-X_{t-1}=(1-B)X_t$ est stationnaire (combinaison finie d'un AR(1)
stationnaire). (A) a une tendance $t$ ; (C) est une marche aléatoire ; (D) a $|\phi^{-1}|>1$
(explosif).
\end{enumerate}

\textbf{Q\,J.2 --- AR(1), observations $-1$ (hier), $+1$ (aujourd'hui).}
\begin{enumerate}[label=\arabic*),leftmargin=*]
\item Prévision demain $=\phi\cdot X_{\text{aujourd'hui}}=\phi\cdot(+1)=$ \bonne{(A) $\phi$}.
\item Variance de l'erreur à 2 jours $=(1+\phi^2)\sigma^2=$ \bonne{(C) $1+\phi^2$} (avec
$\sigma^2=1$).
\end{enumerate}

\textbf{Q\,J.3 --- ARMA(1,1) $X_t=\phi X_{t-1}+\eta_t+\theta\eta_{t-1}$.}
\begin{enumerate}[label=\arabic*),leftmargin=*]
\item Stationnarité : \bonne{$|\phi|<1$} (réponse D). $\theta$ n'intervient que pour
l'inversibilité.
\item Vraie : \bonne{(D)} $\rho(h)=\phi\rho(h-1)$ pour $h$ assez grand (au-delà de $q=1$,
l'ACF d'un ARMA suit la récurrence AR).
\item Si $\phi+\theta=0$ : $\Theta(B)=1+\theta B=1-\phi B=\Phi(B)$, le facteur se simplifie et
$X_t=\eta_t$ : \bonne{(B) un bruit blanc}.
\end{enumerate}

\textbf{Q\,J.4 --- AR$(p)$, affirmation générale vraie.}
\bonne{(B)} : la FAP $P(h)=0$ pour $h>p$. (A) concerne un MA ; (C),(D) ne sont pas générales.

\textbf{Q\,J.5 --- Marche aléatoire, écart-type de $Y_2$.}
$Y_2=\eps_1+\eps_2$, $\Var(Y_2)=2\sigma_\eps^2=2(0.18)=0.36$, écart-type $=\mathbf{0.6}$
(réponse A).

\textbf{Q\,J.6 --- $Y_t=\tfrac12(\eps_t+\eps_{t-1})$, $\sigma_\eps^2=8$.}
$\Cov(Y_t,Y_{t-1})=\tfrac14\Cov(\eps_t+\eps_{t-1},\eps_{t-1}+\eps_{t-2})
=\tfrac14\sigma_\eps^2=\mathbf{2}$.

\textbf{Q\,J.7 --- Moyenne mobile de 8 bruits.}
$\Var(Y_t)=\tfrac{8}{64}\sigma^2$, $\Cov(Y_t,Y_{t-1})=\tfrac{7}{64}\sigma^2$ (7 termes
communs), d'où $\Corr=\tfrac{7}{8}=\mathbf{0.875}$ (indépendant de $\sigma^2$).

\textbf{Q\,J.8 --- AR(1), $\gammahat(0)=0.348$, $\Var(\bar Y)=0.006$, $n=200$.}
Pour un AR(1) : $\Var(\bar Y)\approx\dfrac{\gamma(0)}{n}\cdot\dfrac{1+\phi}{1-\phi}$. D'où
$\dfrac{1+\phi}{1-\phi}=\dfrac{n\,\Var(\bar Y)}{\gamma(0)}=\dfrac{200(0.006)}{0.348}=3.448$,
soit $\phi=\dfrac{3.448-1}{3.448+1}=\mathbf{0.55}$ (réponse B).

\textbf{Q\,J.9 --- MA(2) $\theta_1=0.2,\theta_2=-0.7,\sigma^2=1$.}
$\Var(Y)=1+\theta_1^2+\theta_2^2=\mathbf{1.53}$ ;
$\rho_1=\dfrac{\theta_1+\theta_1\theta_2}{1.53}=\dfrac{0.06}{1.53}=\mathbf{0.039}$ ;
$\rho_2=\dfrac{\theta_2}{1.53}=\mathbf{-0.458}$.

\textbf{Q\,J.10 --- AR(1) $\phi=-0.5,\sigma_\eps=8.66$.}
$\gamma_0=\dfrac{\sigma^2}{1-\phi^2}=\dfrac{75}{0.75}=100$ ($\sigma^2=8.66^2=75$).
$\gamma_1=\phi\gamma_0=\mathbf{-50}$ ; $\gamma_2=\phi^2\gamma_0=\mathbf{25}$.

\textbf{Q\,J.11 --- AR(1) $\phi>0,\sigma_\eps^2=5.333,\gamma(2)=3$.}
$\gamma(2)=\phi^2\gamma_0=\phi^2\dfrac{\sigma^2}{1-\phi^2}=3$ $\Rightarrow$ $\phi=\mathbf{0.6}$ ;
$\gamma_0=\dfrac{5.333}{0.64}=8.33$, $\gamma(1)=\phi\gamma_0=\mathbf{5.0}$.

\textbf{Q\,J.12 --- ARMA(1,1) $\phi=-0.5,\theta=-0.6,\sigma_\eps^2=4$.}
Formules ARMA(1,1) ($X_t=\phi X_{t-1}+\eps_t+\theta\eps_{t-1}$) :
\[
\gamma_0=\sigma^2\,\frac{1+2\phi\theta+\theta^2}{1-\phi^2}
=4\cdot\frac{1+0.6+0.36}{0.75}=\mathbf{10.45} ;\quad
\gamma_1=\sigma^2\,\frac{(1+\phi\theta)(\phi+\theta)}{1-\phi^2}
=4\cdot\frac{(1.3)(-1.1)}{0.75}=\mathbf{-7.63}.
\]
$\rho_1=\gamma_1/\gamma_0=\mathbf{-0.730}$ ; $\rho_2=\phi\rho_1=\mathbf{0.365}$.
\begin{remarque}
Les valeurs proposées dans le QCM d'origine (de l'ordre de $4$--$5$) supposent une convention
ou des paramètres différents ; le calcul ci-dessus est rigoureux pour les paramètres
\emph{donnés}.
\end{remarque}

\textbf{Q\,J.13 --- ARIMA(2,1,2).}
$Y_t=1+1.8Y_{t-1}-1.4Y_{t-2}+0.6Y_{t-3}+\eps_t-0.4\eps_{t-1}+0.1\eps_{t-2}$. La partie AR a
pour polynôme (sur $Y$) $1-1.8B+1.4B^2-0.6B^3$. Comme $d=1$, on factorise $(1-B)$ :
$1-1.8B+1.4B^2-0.6B^3=(1-B)(1-0.8B+0.6B^2)$. La partie AR stationnaire est donc
$1-0.8B+0.6B^2$ : \bonne{$\phi_1=0.8$} (réponse E) et \bonne{$\phi_2=-0.6$} (réponse A).

\textbf{Q\,J.14 --- $Y_t=X_t+\eps_t$, $X_t=X_{t-1}+u_t$.}
$\Delta Y_t=\Delta X_t+\Delta\eps_t=u_t+\eps_t-\eps_{t-1}$.
$\Var(\Delta Y_t)=\sigma_u^2+2\sigma_\eps^2=1+4=\mathbf{5}$ ;
$\Cov(\Delta Y_t,\Delta Y_{t-1})=-\sigma_\eps^2=-2$ ;
$\Corr=\dfrac{-2}{5}=\mathbf{-0.4}$.

\textbf{Q\,J.15 --- AR(2) $\phi_1=-0.5,\phi_2=-0.3$.}
$\rho_1=\dfrac{\phi_1}{1-\phi_2}=\dfrac{-0.5}{1.3}=\mathbf{-0.385}$ ;
$\rho_2=\phi_1\rho_1+\phi_2=-0.5(-0.385)-0.3=\mathbf{-0.108}$ ;
$P_2=\phi_2=\mathbf{-0.3}$ (FAP d'un AR(2) à l'ordre 2).

\textbf{Q\,J.16 --- AR(2), $\gammahat(0)=4,\rhohat(1)=\rhohat(2)=0.5$.}
Yule--Walker : $\begin{pmatrix}1&0.5\\0.5&1\end{pmatrix}\begin{pmatrix}\phi_1\\\phi_2\end{pmatrix}
=\begin{pmatrix}0.5\\0.5\end{pmatrix}\Rightarrow\phi_1=\phi_2=\tfrac13$.
$\hat\sigma_\eps^2=\gamma_0\big(1-\phi_1\rho_1-\phi_2\rho_2\big)=4\big(1-\tfrac13(0.5)-\tfrac13(0.5)\big)
=4\cdot\tfrac23=\mathbf{2.667}$ (réponse A).

\textbf{Q\,J.17 --- AR(2), $\phi_1=0.5,\phi_2=0.3,\mu=5$, $Y_1=5,Y_2=7,Y_3=5$.}
Forme centrée $Y_t-\mu=\phi_1(Y_{t-1}-\mu)+\phi_2(Y_{t-2}-\mu)+a_t$ :
$Y_4=\mu+\phi_1(Y_3-\mu)+\phi_2(Y_2-\mu)=5+0.5(0)+0.3(2)=\mathbf{5.6}$.
\begin{remarque}
La meilleure prévision de $Y_4$ vaut $5.6$. (Les options négatives du QCM d'origine
correspondent vraisemblablement à une autre question ; le calcul est rigoureux.)
\end{remarque}

\textbf{Q\,J.18 --- $Y_t=X_t+V_t$, $X_t$ AR(1).}
$X_t=\phi X_{t-1}+a_t$, $\gamma_X(k)=\dfrac{\sigma_a^2}{1-\phi^2}\phi^{|k|}$. Comme $V_t$ est un
bruit blanc indépendant :
\[
\gamma_Y(0)=\frac{\sigma_a^2}{1-\phi^2}+\sigma_v^2,\qquad
\gamma_Y(k)=\frac{\sigma_a^2}{1-\phi^2}\,\phi^{|k|}\ (k\ne0).
\]
$\gamma_Y(k)$ ne dépend que de $|k|$ et $\E(Y_t)=0$ : \bonne{$Y_t$ est stationnaire (SSL)}.
\begin{remarque}
$Y_t$ est en fait un ARMA(1,1) (somme AR(1) $+$ bruit blanc) : voir Q\,J.19.
\end{remarque}

\textbf{Q\,J.19 --- $y_t=x_t+v_t$, $\phi=0.5,\sigma_v^2=0.25,\sigma_u^2=1$.}
\begin{enumerate}[label=\alph*.]
\item $\gamma_x(k)=\dfrac{\sigma_u^2}{1-\phi^2}\phi^{|k|}=\dfrac{1}{0.75}(0.5)^{|k|}=\tfrac43(0.5)^{|k|}$.
D'où $\gamma_y(0)=\tfrac43+0.25=\dfrac{19}{12}\approx1.583$ et
$\gamma_y(k)=\tfrac43(0.5)^{|k|}$ ($k\ne0$). Indépendant de $t$ : \bonne{stationnaire}.
\item $z_t=y_t-\phi y_{t-1}=(x_t-\phi x_{t-1})+(v_t-\phi v_{t-1})=u_t+v_t-\phi v_{t-1}$.
C'est une somme de bruits blancs aux retards $0$ et $1$ :
$\gamma_z(0)=\sigma_u^2+(1+\phi^2)\sigma_v^2=1+1.25(0.25)=1.3125$ ;
$\gamma_z(1)=-\phi\sigma_v^2=-0.125$ ; $\gamma_z(k)=0$, $|k|\ge2$.
Donc $z_t$ est un \bonne{MA(1)} (ARMA$(0,1)$).
\item On identifie $z_t=e_t-\theta e_{t-1}$ (variance $\sigma_e^2$) :
$\gamma_z(0)=(1+\theta^2)\sigma_e^2=1.3125$ et $\gamma_z(1)=-\theta\sigma_e^2=-0.125$.
Le rapport $\rho_z(1)=\dfrac{-0.125}{1.3125}=-0.0952$. On résout
$\dfrac{-\theta}{1+\theta^2}=-0.0952\Rightarrow\theta^2-10.5\theta+1=0\Rightarrow
\theta\in\{0.0962,\ 10.40\}$. On retient la racine \bonne{inversible} $\theta\approx0.0962$,
d'où $\sigma_e^2=\dfrac{0.125}{\theta}\approx1.30$.
\end{enumerate}

\textbf{Q\,J.20 --- Variance de la moyenne, $x_t=a_t-0.8a_{t-2}$.}
$\gamma(0)=(1+\theta^2)\sigma_a^2=1.64$ ; $\gamma(1)=0$ ; $\gamma(2)=-\theta\sigma_a^2=-0.8$ ;
$\gamma(k)=0$, $|k|\ge3$. Pour $m=\tfrac14\sum_{i=1}^4x_i$ :
\[
\Var(m)=\frac1{16}\sum_{i=1}^4\sum_{j=1}^4\gamma(i-j)
=\frac1{16}\big[4\gamma(0)+2\cdot2\gamma(2)\big]
=\frac1{16}\big[4(1.64)+4(-0.8)\big]=\frac{3.36}{16}=\mathbf{0.21}.
\]
(Les paires d'indices à distance 2 sont $(1,3),(2,4)$ et leurs symétriques, soit
$2\times2=4$ termes $\gamma(2)$.)

\textbf{Q\,J.21 --- Ljung--Box, ARMA(1,1), $K=14$, $\mathrm{LB}=18.386$.}
Sous $H_0$ (résidus $=$ bruit blanc), $\mathrm{LB}\sim\chi^2_{K-p-q}=\chi^2_{14-1-1}=\chi^2_{12}$.
Valeur critique à $5\%$ : $\chi^2_{12,0.95}=21.03$. Comme $18.386<21.03$,
\begin{reponse}on \bonne{ne rejette pas $H_0$} : les résidus sont compatibles avec un bruit
blanc, le modèle ARMA(1,1) est \emph{validé} (réponse : \og on ne rejette pas $H_0$ \fg).
\end{reponse}

% =====================================================================
\section{Hors programme du cours BADAOUI\,\texorpdfstring{\horsprog}{*}}
\markboth{Hors programme}{}
\begin{tcolorbox}[colback=orangelight,colframe=orange,title={Notions au-delà du cours BADAOUI}]
Ces corrigés mobilisent des notions (bruit blanc fort/faible, stationnarité mutuelle,
projection orthogonale dans $L^2$, inversion formelle d'opérateurs) qui \emph{dépassent} le
cours étudié. Ils sont fournis pour complétude.
\end{tcolorbox}

\textbf{Q\,Z.1 --- $\eps_t=\eta_t\eta_{t-1}\eta_{t-2}$.}
$\E(\eps_t)=\E(\eta_t)\E(\eta_{t-1})\E(\eta_{t-2})=0$ ; $\Var(\eps_t)=\E(\eta_t^2)^3=1$ ;
pour $k\ne0$, $\E(\eps_t\eps_{t-k})=0$ (un facteur $\eta$ apparaît à une puissance impaire,
d'espérance nulle, par indépendance). Donc $(\eps_t)$ est un \bonne{bruit blanc faible}. Il
n'est \emph{pas} fort (les $\eps_t$ ne sont pas indépendants : $\eps_t$ et $\eps_{t-1}$
partagent $\eta_{t-1},\eta_{t-2}$). Enfin $\eps_t^2=\eta_t^2\eta_{t-1}^2\eta_{t-2}^2$ est
stationnaire et son ACF s'annule au-delà du retard 2 : $(\eps_t^2)$ suit un \bonne{MA(2)}.

\textbf{Q\,Z.2 --- Stationnarité mutuelle.}
Deux processus $X,Y$ sont \emph{conjointement} (mutuellement) stationnaires ssi : $X$ est
stationnaire, $Y$ est stationnaire \emph{et} la covariance croisée
$\gamma_{XY}(t,s)=\Cov(X_t,Y_s)$ ne dépend que de $t-s$ (invariante par translation).
\begin{reponse}Cocher : \bonne{$X$ stationnaire}, \bonne{$Y$ stationnaire}, \bonne{covariance
croisée invariante par translation}.\end{reponse}

\textbf{Q\,Z.3 --- Inverses d'opérateurs.}
$(I-\tfrac15B)^{-1}=\sum_{k\ge0}\big(\tfrac15\big)^kB^k$ (converge car $|\tfrac15|<1$).
$(I+3B)^{-1}$ : la série $\sum(-3)^kB^k$ \emph{diverge} ($|3|>1$) ; on inverse \og vers le
futur \fg{} : $(I+3B)^{-1}=-\tfrac13B^{-1}\big(I+\tfrac13B^{-1}\big)^{-1}
=-\sum_{k\ge1}(-\tfrac13)^kB^{-k}$ (développement en puissances de $B^{-1}$).

\textbf{Q\,Z.4 --- Projection orthogonale, $\Phi(B)=I-\tfrac16B-\tfrac16B^2$.}
\begin{enumerate}[label=\arabic*),leftmargin=*]
\item Racines de $1-\tfrac16z-\tfrac16z^2=0$ : $z=\dfrac{-1\pm\sqrt{1+24}}{2}\cdot(-1)\cdots$
on résout $z^2+z-6=0$ (en multipliant par $-6$) $\Rightarrow z\in\{2,-3\}$, modules $>1$ :
\bonne{$\Phi$ inversible}. $\Phi(B)^{-1}=\sum_{k\ge0}c_kB^k$ avec $c_0=1$,
$c_k=\tfrac16c_{k-1}+\tfrac16c_{k-2}$.
\item $\Phi(B)X_t=\eps_t$ : \bonne{AR(2)} causal stationnaire, $p=2,q=0$ ;
$X_t=\Phi(B)^{-1}\eps_t=\sum_{k\ge0}c_k\eps_{t-k}$ (MA$(\infty)$) ; $\E(X_t)=0$ ;
$\gamma(h)=\sigma^2\sum_{k\ge0}c_kc_{k+|h|}$ ; dans le cas gaussien $(X_t,X_{t+3})$ suit une loi
normale bivariée $\mathcal N\!\big(0,\,\begin{psmallmatrix}\gamma(0)&\gamma(3)\\\gamma(3)&\gamma(0)\end{psmallmatrix}\big)$.
\item $\Phi(B)Y_t=\eps_t-\tfrac13\eps_{t-1}$ : \bonne{ARMA(2,1)} ; stationnaire (mêmes racines
AR) et inversible ($1-\tfrac13B$, racine $3$). La représentation est \emph{minimale} car
$\Theta(B)=1-\tfrac13B$ ne partage aucune racine avec $\Phi(B)$ (racines $2,-3$). $(\eps_t)$
\emph{est} le processus des innovations (MA inversible). La forme AR$(\infty)$ s'obtient par
$\eps_t=\Theta(B)^{-1}\Phi(B)Y_t$ ; la projection orthogonale de $Y_t$ sur le passé
$\overline{\text{vect}}(Y_{t-1},Y_{t-2},\dots)$ est $\widehat Y_t=Y_t-\eps_t$ (résidu
d'innovation), soit la partie \og prévisible \fg{} de $Y_t$.
\end{enumerate}

\textbf{Q\,Z.5 --- Lissage et filtres (QCM).}
\begin{enumerate}[label=\arabic*),leftmargin=*]
\item \bonne{Vrai} (le lissage exponentiel simple ajuste localement une constante).
\item \bonne{Vrai} (Holt--Winters intègre une composante saisonnière).
\item Pour estimer la tendance $m$ en présence de saisonnalité $s$ : filtre qui
\bonne{absorbe $s$ et laisse invariante $m$} (une MM d'ordre égal à la période).
\item Pour tester la blancheur des résidus : \bonne{regarder ACF/PACF empiriques} et
\bonne{utiliser un test de Ljung--Box} (l'AIC sert à la sélection, pas au test de blancheur).
\end{enumerate}

\vfill
\begin{center}\small\color{gris}
Fin du corrigé --- document fondé sur le cours de Séries Chronologiques de F.~BADAOUI (INSEA).
\end{center}

\end{document}
